\documentclass[12pt]{article}
\usepackage[utf8]{inputenc}
\usepackage[T1]{fontenc}
\usepackage[margin=1in]{geometry}
\usepackage{setspace}
\usepackage{xcolor}
\usepackage[colorlinks=true, urlcolor=blue!60!black, linkcolor=blue!60!black]{hyperref}
\setlength{\parindent}{0em}
\setlength{\parskip}{0.6em}
\renewcommand{\baselinestretch}{1.0}
\begin{document}
\begin{center}
{\LARGE \textit{R.A.V. v. City of St. Paul}, 505 U.S. 377 (1992)\par}
\vspace{0.5em}
{\large Case Brief --- Two-Page Edition\par}
\end{center}

\section{Memory Jogger}
St.~Paul's cross-burning ordinance fell because, even within unprotected fighting words, government may not single out disfavored topics or viewpoints (\href{https://supreme.justia.com/cases/federal/us/505/377/}{505 U.S. 377}).

\section{Detailed Case Facts}
In 1990, R.A.V., a juvenile, allegedly burned a makeshift cross in the fenced yard of a Black family in St.~Paul. Rather than use available content-neutral charges, the City charged him under an ordinance banning symbols, including a burning cross or swastika, that one knows arouse anger, alarm, or resentment on the basis of race, color, creed, religion, or gender.

\section{Procedural History}
The trial court dismissed; the Minnesota Supreme Court reversed, limiting the ordinance to fighting words; the Supreme Court reversed.

\section{Judicial Votes}
9--0 in judgment, 5--4 on rationale. Scalia wrote for Rehnquist, Kennedy, Souter, and Thomas; White, Blackmun, and Stevens concurred in the judgment.

\section{Holding}
Even limited to fighting words, the ordinance is facially invalid content and viewpoint discrimination within a proscribable category, unjustified by strict scrutiny.

\section{Analysis of Opinions}
Scalia held that proscribable categories are not invisible to the Constitution: they may be regulated for their proscribable content, not their message. Exceptions cover distinctions tracking the reason for proscription, secondary effects or conduct, and cases with no realistic risk of idea suppression. The ordinance singled out five topics and, in practice, one viewpoint, and a neutral ban on all fighting words would serve the City's interest. White would have held the ordinance overbroad and attacked the majority's underbreadth theory; Stevens favored a contextual test.

\section{Examples: Future Applications}
\textbf{Same-Side:} a ban only on fighting words disparaging police is invalid. \textbf{Opposite-Side:} enhanced penalties for threats against election workers fit the first exception. \textbf{Fence-Sitter:} a campus ban on fighting words demeaning nationality, religion, or political ideology.

\section{Critique}
Critics see a paradox in preferring broader bans and an elastic first exception, as \textit{Virginia v.\ Black} shows; defenders say \textit{R.A.V.} stops majorities from silencing disfavored views.

\section{Key Quotations}
\begin{quote}
``What they mean is that these areas of speech can, consistently with the First Amendment, be regulated because of their constitutionally proscribable content (obscenity, defamation, etc.)---not that they are categories of speech entirely invisible to the Constitution, so that they may be made the vehicles for content discrimination unrelated to their distinctively proscribable content.''
\end{quote}
\noindent --- Justice Scalia, Opinion of the Court, 505~U.S.~at~383--384.

\end{document}
